\section*{Bab \ref{ch:b1:fluid-envelopes-structure} --- Struktur atmosfer dan samudra}
% values: calc/ch26.py

\begin{solution}{exo:b1:fluid-envelopes-structure:1}
(a)~$H = R^*T/(M_0 g_0) = 8314.32 \times 250/(28.9644 \times 9.80665) =
\qty{7.32}{km}$. (b)~$e^{10/7.32} = 3.92$: tekanan turun hingga sekitar seperempatnya.
(c)~$H$ sebanding dengan suhu: pada \qty{216.65}{K} ia adalah
\qty{6.34}{km}, dibanding \qty{8.43}{km} pada \qty{288.15}{K}. Udara dingin lebih padat
pada tekanan tertentu, sehingga tekanan turun lebih cepat terhadap ketinggian.
\end{solution}

\begin{solution}{exo:b1:fluid-envelopes-structure:2}
Dengan $\Gamma_d = \qty{9.76}{K/km}$: \qtyrange{0}{0.3}{km}, $11.0 > 9.76$,
tidak stabil; \qtyrange{0.3}{1.2}{km}, $9.7$, netral dalam ketelitian data;
\qtyrange{1.2}{1.5}{km}, inversi; \qtyrange{1.5}{8}{km}, stabil;
di atas \qty{8}{km}, isotermal, stabil. Inversi adalah yang paling stabil, karena
$\Gamma_d - \Gamma = 13.8$ paling besar di sana.
\end{solution}

\begin{solution}{exo:b1:fluid-envelopes-structure:3}
(a)~$\sigma = 1026.95 - 1000 = \qty{26.95}{kg.m^{-3}}$. (b)~B, yang lebih ringan, berada
di atas A. (c)~Semua massa jenis air laut dimulai dengan digit yang sama, 10; $\sigma$ hanya
menyimpan digit yang berubah dari satu air ke air lain.
\end{solution}

\begin{solution}{exo:b1:fluid-envelopes-structure:4}
(a)~$S = 0.25 \times 35.40 + 0.75 \times 34.30 = 34.575$; $\theta = 0.25 \times
15.0 + 0.75 \times 4.0 = \qty{6.75}{\celsius}$. (b)~Pada ruas PQ, seperempat jalan dari
Q ke P (\cref{prop:b1:fluid-envelopes-structure:mixing}).
(c)~Campuran P, Q dan R mengisi segitiga PQR. Menulis
$(34.60, 8.0) = a\,\mathrm P + b\,\mathrm Q + c\,\mathrm R$ dengan $a + b + c = 1$
memberi $a = 0.345$, $b = 0.788$, $c = -0.133$: proporsi R negatif, sehingga titik itu
berada di luar segitiga dan tidak ada campuran dari ketiganya yang dapat menghasilkannya.
\end{solution}

\begin{solution}{exo:b1:fluid-envelopes-structure:5}
(a)~$\theta = 253.15 \times 2^{2/7} = \qty{308.6}{K}$. (b)~Pada \qty{1000}{hPa},
$\theta = T = \qty{298.15}{K}$. \omterm{def:b1:fluid-envelopes-structure:potential-temperature}{Suhu potensial} naik ke atas, dari 298.2 menjadi
\qty{308.6}{K}: lapisan itu stabil untuk udara kering. (c)~Ia akan tiba di
\omterm{def:b1:fluid-envelopes-structure:potential-temperature}{suhu potensialnya}, \qty{308.6}{K}, yaitu \qty{35.4}{\celsius}, sepuluh derajat lebih hangat daripada udara yang sudah ada di sana.
\end{solution}

\begin{solution}{exo:b1:fluid-envelopes-structure:6}
(a)~$T = 288.15 - 6.5 \times 5 = \qty{255.65}{K}$ dan $p = \num{101325} \times
(255.65/288.15)^{5.256} = \qty{54020}{Pa}$, \qty{540.2}{hPa}. (b)~$H =
\qty{8.43}{km}$ dan $p = \num{101325}\,e^{-5/8.43} = \qty{56010}{Pa}$. Tekanan
yang sebenarnya (standar) lebih rendah: udara di sana lebih dingin daripada \qty{288.15}{K}, sehingga
lebih padat, dan tekanan turun lebih cepat terhadap ketinggian. (c)~$p = \num{22632}\,
e^{-9/6.34} = \qty{5475}{Pa}$, nilai dari
\cref{tab:b1:fluid-envelopes-structure:layers}.
\end{solution}

\begin{solution}{exo:b1:fluid-envelopes-structure:7}
(a)~$\beta\,\Delta S = \alpha\,\Delta T$ memberi $\Delta S = 1.67/7.61 = 0.22$.
(b)~$\Delta\rho = 1026.95 \times (-1.67\times10^{-4} \times 2 + 7.61\times10^{-4}
\times 0.3) = \qty{-0.11}{kg.m^{-3}}$: air itu menjadi lebih ringan.
(c)~$\Delta\rho = 1028.11 \times (-0.53\times10^{-4} \times 2 + 7.85\times10^{-4}
\times 0.3) = \qty{+0.13}{kg.m^{-3}}$: perubahan yang sama kini membuatnya lebih padat, karena
pemanasan hampir tidak mengubah massa jenis air dekat titik beku.
\end{solution}

\begin{solution}{exo:b1:fluid-envelopes-structure:8}
(a)~$N^2 = (9.81/260) \times (9.76 - 4)\times10^{-3} = \qty{2.17e-4}{s^{-2}}$,
$N = \qty{0.0147}{s^{-1}}$, periode $2\pi/N = \qty{426}{s}$, sekitar tujuh
menit. (b)~$N^2 = (9.81/270) \times 14.76\times10^{-3} = \qty{5.36e-4}{s^{-2}}$,
$N = \qty{0.0232}{s^{-1}}$, periode \qty{271}{s}. (c)~$N$ mendekati nol dan periode
tumbuh tanpa batas: gaya pemulih lenyap pada lapisan netral.
\end{solution}

\begin{solution}{exo:b1:fluid-envelopes-structure:9}
(a)~Ia lebih padat, $\sigma_\theta = 27.20$ dibanding 26.30: AAIW lebih tawar tetapi
jauh lebih dingin. (b)~$\Delta\sigma = 1026.95 \times [-1.67\times10^{-4} \times
(4.31 - 15.43) + 7.61\times10^{-4} \times (34.30 - 35.54)] = 1.907 - 0.969 =
0.94$, dibanding $0.90$ secara eksak. (c)~Suhu: pendinginan menambah
\qty{1.91}{kg.m^{-3}}, pengurangan salinitas menghilangkan \qty{0.97}{kg.m^{-3}}.
\end{solution}

\begin{solution}{exo:b1:fluid-envelopes-structure:10}
(a)~Dengan \cref{thm:b1:fluid-envelopes-structure:stability}, $(g/\theta)\,\mathrm
d\theta/\mathrm dz = (g/T)(\Gamma_d - \Gamma)$; $\mathrm d\theta/\mathrm dz = 0$
memberi $\Gamma = \Gamma_d$. (b)~$T = T_s - \Gamma_d z$, lapisan dengan gradien konstan
$L = \Gamma_d$: dengan \cref{prop:b1:fluid-envelopes-structure:layer-pressure}, $p =
p_s (T/T_s)^{gM/(R\Gamma_d)}$, dan $gM/(R\Gamma_d) = gMc_p/(Rg) = c_pM/R$
($= 3.5$ untuk udara kering). (c)~$p$ lenyap di mana $T = 0$, pada $z = T_s/\Gamma_d =
300/9.76 = \qty{30.7}{km}$. Lapisan nyata dengan \omterm{def:b1:fluid-envelopes-structure:potential-temperature}{suhu potensial} seragam paling
dalam hanya beberapa kilometer: mereka ditutup oleh udara stabil, seperti troposfer
oleh stratosfer.
\end{solution}

\begin{solution}{exo:b1:fluid-envelopes-structure:11}
(a)~$S = (33.787 + 36.050)/2 = 34.919$, $\theta = \qty{8.0}{\celsius}$.
(b)~Dengan persamaan keadaan linear, $\sigma$ adalah fungsi affin dari $S$ dan
$\theta$, sehingga $\sigma$ campuran adalah rerata dari keduanya, 27.00.
(c)~$\alpha$ adalah \qty{0.75e-4}{K^{-1}} pada \qty{2}{\celsius} dan
\qty{2.07e-4}{K^{-1}} pada \qty{14}{\celsius}, sedangkan $\beta$ sedikit berubah
(\num{7.80e-4} dan \num{7.53e-4}). Massa jenis turun semakin cepat seiring naiknya suhu:
ia adalah fungsi cekung terhadap suhu. Sepanjang garis pencampuran, campuran berada pada
suhu rata-rata, tempat massa jenisnya berada di atas tali yang menghubungkan kedua induknya,
yaitu di atas 27.00. (d)~Kelebihannya berasal dari lengkungan, suku orde dua: ia tumbuh
sebagai kuadrat selisih suhu kedua induknya.
\end{solution}

\begin{solution}{exo:b1:fluid-envelopes-structure:12}
(a)~Selisih tekanannya $5604 - 5500 = \qty{104}{dbar}$: parsel mendingin
$0.129 \times 0.104 = \qty{0.013}{\celsius}$, menjadi \qty{0.718}{\celsius}.
(b)~Air di sekitarnya pada \qty{5400}{m} bersuhu \qty{0.730}{\celsius}, lebih hangat: parsel,
dengan salinitas yang hampir sama, lebih padat daripada sekitarnya dan tenggelam kembali. Lapisan itu stabil, meskipun suhu terukur naik ke bawah; \omterm{def:b1:fluid-envelopes-structure:potential-temperature}{suhu potensialnya} turun ke bawah.
(c)~$\theta \approx
0.7313 - 0.086 \times 5.604 = \qty{0.249}{\celsius}$, dalam seperseribu derajat dari nilai eksak.
\end{solution}

\begin{solution}{pb:b1:fluid-envelopes-structure:1}
\begin{enumerate}
  \item Suhu hampir seragam hingga \qty{50}{m}, turun cepat hingga
        \qty{300}{m}, lebih lambat hingga \qty{1000}{m}, dan hampir seragam di bawahnya.
        Salinitas turun dari 35.80 di permukaan hingga minimum 34.32 pada
        \qty{800}{m}, lalu naik hingga 34.90 pada \qty{3000}{m}.
  \item $22.40 - 0.5 = \qty{21.90}{\celsius}$ tercapai di antara \qty{50}{m}
        (22.10) dan \qty{75}{m} (19.20): $50 + 25 \times 0.20/2.90 = \qty{52}{m}$.
  \item $(22.10 - 13.00)/250 = \qty{0.0364}{\celsius/m}$, \qty{3.6}{\celsius}
        per \qty{100}{m}. \omterm{def:b1:fluid-envelopes-structure:ocean-layers}{Termoklin} membentang dari dasar \omterm{def:b1:fluid-envelopes-structure:ocean-layers}{lapisan campuran}, dekat
        \qty{50}{m}, hingga sekitar \qty{800}{m}; ia paling tajam antara
        \qty{50}{m} dan \qty{75}{m} (\qty{11.6}{\celsius} per \qty{100}{m}).
  \item $\sigma = 1026.95 \times [1 - 1.67\times10^{-4} \times 12.4 +
        7.61\times10^{-4} \times 0.80] - 1000 = 25.45$, \qty{0.74}{kg.m^{-3}} terlalu
        tinggi. Air permukaan jauh dari keadaan acuan, dan $\alpha$ jauh lebih besar pada
        \qty{22}{\celsius} daripada pada \qty{10}{\celsius}: bentuk linear meremehkan
        keringanan akibat kehangatan.
  \item Minimum salinitas pada \qty{800}{m}, 34.32 pada \qty{4.8}{\celsius}, adalah
        Air Menengah Antartika, terbentuk di bawah angin barat yang hujan di Samudra
        Selatan. Maksimum pada \qty{3000}{m}, 34.90 pada
        \qty{2.5}{\celsius}, adalah Air Dalam Atlantik Utara, terbentuk di Laut
        Greenland dan Labrador.
  \item $(2.50 - 2.26)/3.042 = \qty{0.079}{\celsius}$ per \qty{1000}{dbar}.
  \item $\sigma_\theta$ bertambah dengan kedalaman: kolom itu stabil. \omterm{def:b1:fluid-envelopes-structure:ocean-layers}{Piknoklin}
        berada di tempat $\sigma_\theta$ bertambah paling cepat, di beberapa ratus meter
        teratas, bersama \omterm{def:b1:fluid-envelopes-structure:ocean-layers}{termoklin}: \num{0.82} antara \qty{200}{m} dan \qty{800}{m},
        hanya \num{0.70} sepanjang \qty{2200}{m} di bawahnya.
  \item $S = (35.45 + 34.32)/2 = 34.885$; $\theta = (14.90 + 4.80)/2 =
        \qty{9.85}{\celsius}$.
  \item C: 26.461; I: 27.310; campuran: 26.886, tepat rerata dari keduanya
        (26.886), karena bentuk linear bersifat affin.
  \item Rerata induknya adalah $(26.342 + 27.159)/2 = 26.751$; campuran,
        pada 26.888, \qty{0.14}{kg.m^{-3}} lebih padat. Ini adalah cabbeling.
  \item Dengan bentuk linear, pengukuran memberi 26.61 pada \qty{300}{m}, 26.88 pada
        \qty{500}{m} dan 27.31 pada \qty{800}{m}: massa jenis campuran berada dekat
        \qty{500}{m}. Air di sana, $(34.82, \qty{9.6}{\celsius})$, sangat dekat dengan
        campuran $(34.885, \qty{9.85}{\celsius})$: pengukuran di antara Air Sentral dan
        AAIW sendiri hampir merupakan garis pencampuran di antara keduanya.
  \item $S = 0.25 \times 35.45 + 0.75 \times 34.32 = 34.60$, $\theta = 0.25
        \times 14.90 + 0.75 \times 4.80 = \qty{7.33}{\celsius}$. Interpolasi profil
        antara \qty{500}{m} dan \qty{800}{m} memberi
        $(34.57, \qty{7.2}{\celsius})$ dekat \qty{650}{m}: hampir tepat berada di situ.
  \item Dalam \si{K/km}: A--B 10.0; B--C 9.7; C--D $-5.5$; D--E 6.29; E--F 6.50;
        F--G 6.57; G--H 0.33; H--I $-0.50$; I--J $-0.83$.
  \item A--B tidak stabil; B--C netral; C--D inversi; D--E, E--F, F--G stabil;
        G--H hampir isotermal, stabil; H--I dan I--J inversi. Tropopause
        berada di G, \qty{12.5}{km} dan \qty{-59.0}{\celsius}, tempat suhu
        berhenti menurun.
  \item Suhu rata-rata adalah \qty{214.65}{K}: $H = 287 \times 214.65/9.81 =
        \qty{6.28}{km}$, dan $147\,e^{-4/6.28} = \qty{77.7}{hPa}$, konsisten
        dengan 78 yang terukur.
  \item Udara di A, \qty{22.0}{\celsius}, lebih dingin daripada laut: panas mengalir
        dari laut ke udara, yang dipanaskan dari bawah. Lapisan A--B, dengan laju lapse
        di atas $\Gamma_d$, adalah hasilnya.
  \item Di B: $22.0 - 9.76 \times 0.5 = \qty{17.1}{\celsius}$, dibanding 17.0; di
        C: \qty{7.4}{\celsius}, dibanding 7.3; di D: \qty{3.5}{\celsius}, dibanding
        9.5. Parsel sedikit lebih hangat daripada udara hingga C, dan jauh lebih dingin di D.
  \item Di dasar inversi, dekat \qty{1.5}{km}: saat memasuki inversi, parsel
        menjadi lebih dingin, dan karenanya lebih padat, daripada udara di sekitarnya,
        sehingga kenaikannya terhenti.
  \item Parsel yang dilepaskan oleh lapisan permukaan yang tidak stabil naik melewatinya
        dan mencampurnya hingga \omterm{def:b1:fluid-envelopes-structure:potential-temperature}{suhu potensialnya} seragam, sehingga laju lapsenya
        menjadi $\Gamma_d$; inversi di atasnya menghentikan pencampuran.
  \item Dalam K: A 294.15; B 293.99; C 294.08; D 300.51; E 312.65; F 329.29; G
        346.28; H 369.50; I 446.99; J 600.91.
  \item $\theta$ menurun dari A ke B (tidak stabil), hampir konstan dari B ke
        C (netral) dan naik di atasnya (stabil), paling cepat dari C ke D dan di atas
        G: klasifikasi yang sama seperti pada soal~14.
  \item E--F: $(329.29 - 312.65)/4 = \qty{4.2}{K/km}$; H--I: $(446.99 -
        369.50)/4 = \qty{19.4}{K/km}$, hampir lima kali lebih besar.
  \item Dengan rerata $\theta = \qty{408.2}{K}$: $N^2 = (9.81/408.2) \times
        0.01937 = \qty{4.66e-4}{s^{-2}}$, $N = \qty{0.0216}{s^{-1}}$, periode
        \qty{291}{s}, sekitar lima menit.
  \item $N^2 = (9.81/1027) \times (26.35 - 24.79)/150 = \qty{9.9e-5}{s^{-2}}$,
        $N = \qty{0.0100}{s^{-1}}$, periode \qty{630}{s}. Stratosfer bawah lebih stabil,
        dengan $N^2$ hampir lima kali lebih besar daripada di \omterm{def:b1:fluid-envelopes-structure:ocean-layers}{termoklin}.
  \item Di G, $\theta = 214.15 \times (1000/186)^{2/7} = \qty{346}{K}$: \omterm{def:b1:fluid-envelopes-structure:potential-temperature}{suhu
        potensial} pada tropopause sekitar \qty{346}{K}.
\end{enumerate}
\end{solution}
